\documentclass[12pt,a4paper]{article}
\usepackage[a4paper,left=25mm,right=20mm,top=24mm,bottom=24mm]{geometry}
\usepackage{fontspec}
\usepackage{polyglossia}
\setdefaultlanguage{russian}
\setotherlanguage{english}
\IfFontExistsTF{Libertinus Serif}{
    \setmainfont{Libertinus Serif}
    \newfontfamily\cyrillicfont{Libertinus Serif}
}{
    \IfFontExistsTF{DejaVu Serif}{
        \setmainfont{DejaVu Serif}
        \newfontfamily\cyrillicfont{DejaVu Serif}
    }{
        \setmainfont{Times New Roman}
        \newfontfamily\cyrillicfont{Times New Roman}
    }
}
\IfFontExistsTF{DejaVu Sans Mono}{\setmonofont{DejaVu Sans Mono}[Scale=0.82]}{}
\usepackage{microtype}
\usepackage{setspace}
\setstretch{1.20}
\usepackage{parskip}
\setlength{\parindent}{0pt}
\usepackage{booktabs}
\usepackage{tabularx}
\usepackage{longtable}
\usepackage{array}
\usepackage{ragged2e}
\newcolumntype{P}[1]{>{\RaggedRight\arraybackslash}p{#1}}
\newcolumntype{Y}{>{\RaggedRight\arraybackslash}X}
\setlength{\emergencystretch}{3em}
\usepackage{enumitem}
\usepackage{needspace}
\usepackage{amsmath}
\usepackage{amssymb}
\setlist{topsep=3pt,itemsep=2pt,parsep=0pt}
\usepackage{graphicx}
\usepackage{xcolor}
\definecolor{codegray}{RGB}{62,66,72}
\usepackage{listings}
\lstset{
    basicstyle=\ttfamily\footnotesize,
    breaklines=true,
    breakatwhitespace=false,
    columns=fullflexible,
    showstringspaces=false,
    frame=single,
    rulecolor=\color{black!20},
    xleftmargin=4pt,
    xrightmargin=4pt,
    aboveskip=8pt,
    belowskip=8pt
}
\usepackage{tikz}
\usetikzlibrary{arrows.meta,positioning,shapes.geometric}
\usepackage{hyperref}
\hypersetup{hidelinks,pdfauthor={Agibalov Vladislav},pdflang={ru-RU}}
\usepackage{fancyhdr}
\pagestyle{fancy}
\setlength{\headheight}{14pt}
\fancyhf{}
\fancyhead[L]{\small mangen}
\fancyhead[R]{\small Техническая документация · 2026}
\fancyfoot[C]{\thepage}
\renewcommand{\headrulewidth}{0.2pt}
\setcounter{tocdepth}{2}
\newcommand{\impl}{\textbf{Реализация.}~}

\title{mangen\\Руководство пользователя и сборки}
\author{Агибалов Владислав Сергеевич\\НИЯУ МИФИ, ИВТ}
\date{2026 год}

\begin{document}
\maketitle
\begin{abstract}
Руководство по автономной сборке, запуску, интерпретации результатов,
проверке хэш-сумм и аутентичности манифеста. Описываются различия между
macOS и Linux, ограничения доступа к файлам и двухкоммитный выпуск ZIP.
\end{abstract}
\tableofcontents
\newpage

\section{Назначение и область применения}

\texttt{mangen} предназначен для инвентаризации содержимого файлового
каталога с криптографическими хэшами SHA-2. Программа не изменяет файлы
обрабатываемого дерева и не создаёт манифест самостоятельно: перенаправление
\texttt{stdout} выполняется оболочкой. Для записи каталога в манифест
необходимо исключить его из обхода или хранить вывод вне проверяемого дерева.

\section{Подготовка к автономной сборке}

Требуются CMake версии, заявленной в \texttt{CMakeLists.txt},
компилятор C17, POSIX-совместимая система macOS или Linux и инструменты сборки
(Make либо Ninja). Для \texttt{BUILD\_TESTING=ON} дополнительно нужен
компилятор C++17. Библиотеки Concurrency Kit, Mbed TLS и GoogleTest должны
находиться в \texttt{libs/third\_party/} до начала конфигурации.

\subsection{Сборка с тестами}

Команды выполняются в корне распакованного исходного каталога:

\begin{lstlisting}
cmake -S . -B build -DCMAKE_BUILD_TYPE=Release -DBUILD_TESTING=ON
cmake --build build --parallel 4
ctest --test-dir build --output-on-failure
cmake --install build --prefix "$PWD/build/install"
./build/install/bin/mangen --help
\end{lstlisting}

Выполняется локальная конфигурация; скачивание библиотек во время build не
требуется. Каталог установки задаётся относительным к пользователю префиксом,
поэтому привилегии администратора обычно не нужны.

\subsection{Сборка только производственного приложения}

\begin{lstlisting}
cmake -S . -B build/release -DCMAKE_BUILD_TYPE=Release \
  -DBUILD_TESTING=OFF
cmake --build build/release --parallel 4
cmake --install build/release --prefix "$PWD/build/release/install"
\end{lstlisting}

При этом проект не должен требовать C++ и GoogleTest. Если сборка с
\texttt{BUILD\_TESTING=OFF} требует C++-компилятор, это нарушение контракта
сборки, подлежащее исправлению.

\subsection{macOS и Linux}

\begin{tabularx}{\textwidth}{P{0.28\textwidth}YY}
\toprule
Задача & macOS & Linux \\
\midrule
Компилятор & Apple Clang или GCC & GCC или Clang \\
SHA-256 & \texttt{shasum -a 256} & \texttt{sha256sum} \\
SHA-384 & \texttt{shasum -a 384} & \texttt{sha384sum} \\
SHA-512 & \texttt{shasum -a 512} & \texttt{sha512sum} \\
Диагностика памяти & ASan/UBSan, \texttt{leaks} & ASan/UBSan/LSan \\
\bottomrule
\end{tabularx}

Набор Bash-скриптов использует Bash 3.2 или новее и определяет доступную
системную утилиту SHA без установки сторонних пакетов. Для ОС Linux
действительно требуется отдельный запуск компиляции и CTest, одной
успешной macOS-сборки недостаточно.

\section{Командная строка}

\begin{lstlisting}
mangen [DIR_PATH] [OPTIONS]
\end{lstlisting}

\begin{longtable}{P{0.28\textwidth}P{0.62\textwidth}}
\toprule
Аргумент & Семантика \\
\midrule
\endhead
\texttt{DIR\_PATH} & Корневой каталог, по умолчанию текущий. \\
\texttt{-h}, \texttt{--help} & Справка без сканирования дерева. \\
\texttt{-v}, \texttt{--version} & Информация о версии приложения и сборке. \\
\texttt{-e PATTERN} & Исключение базовых имён, многократно, объединение ИЛИ. \\
\texttt{--sha256} & SHA-256 по умолчанию. \\
\texttt{--sha384} & SHA-384. \\
\texttt{--sha512} & SHA-512. \\
\texttt{-j N}, \texttt{--threads N} & Положительное количество потоков. \\
\texttt{--queue-size N} & Стартовая ёмкость очереди каталогов. \\
\bottomrule
\end{longtable}

\section{Получение и интерпретация манифеста}

\begin{lstlisting}
./build/install/bin/mangen ./data --sha256 --threads 4 \
    > ./manifest.txt 2> ./diagnostics.txt
printf 'exit=%s\n' "$?"
\end{lstlisting}

Одна успешная запись содержит относительный путь, пробел, двоеточие, пробел
и полное HEX-значение SHA. Имена отсортированы. Содержимое не зависит от
порядка завершения рабочих потоков. Ошибки чтения файлов, доступа к каталогам
и некорректных ссылок помещаются в \texttt{diagnostics.txt}.

\subsection{Коды завершения}

\begin{itemize}
\item \texttt{0}: все применимые файлы обработаны без ошибок.
\item \texttt{1}: ошибка во время обхода или хэширования; доступные записи
могут присутствовать, но манифест является \emph{частичным}.
\item \texttt{2}: ошибка аргументов запуска.
\end{itemize}

Нельзя считать частичный манифест полным только потому, что файл
\texttt{manifest.txt} существует и содержит корректные SHA.

\section{Проверка прав на двух деревьях}

Скрипт создаёт отдельные \texttt{clean} и \texttt{restricted} с одинаковым
содержимым. Во втором дереве рандомизируются права chmod; гарантированно
закрываются хотя бы один каталог и один файл. Тест должен выполняться без
root, иначе проверки прав могут оказаться неинформативны.

\begin{lstlisting}
TEST_ROOT=$(bash scripts/generate_test_trees.sh)
bash scripts/verify_test_trees.sh "$TEST_ROOT" \
    ./build/install/bin/mangen
\end{lstlisting}

Проверка создаёт независимый эталон посредством \texttt{find},
\texttt{shasum} или GNU \texttt{sha*sum}, сравнивает полный список строк
побайтно, сортировку и идентичность результатов для одного и четырёх
потоков. Проверяются все три SHA-алгоритма. Статус ограниченного дерева
должен быть \texttt{1} у обычного пользователя; доступная часть
обязана совпадать с эталоном.

\section{Контроль целостности при передаче}

\subsection{Три разных гарантии}

\begin{enumerate}
\item \textbf{SHA-256 ZIP} выявляет случайные изменения байтов архива,
если значение получено из доверенного источника. Hash-файл из того же
недоверенного канала сам по себе не доказывает происхождение.
\item \textbf{Git commit hash} идентифицирует исходное зафиксированное
дерево, но не является суммой ZIP и не удостоверяет несохранённые правки.
\item \textbf{Цифровая SSH-подпись манифеста} подтверждает подлинность
самого документа при условии, что доверенный публичный ключ получен
отдельно и закреплён получателем.
\end{enumerate}

\subsection{Подпись манифеста через OpenSSH}

Для подписания подходит SSH-ключ Ed25519. Команда
\texttt{ssh-keygen} в режиме \texttt{-Y sign} формирует подпись.
Секретный ключ не помещают в ZIP.

\begin{lstlisting}
bash scripts/sign_manifest.sh manifest.txt "$HOME/.ssh/id_ed25519"
\end{lstlisting}

Создаётся \texttt{manifest.txt.sig}. Получатель заранее получает
\texttt{trusted\_signers} через отдельный доверенный канал:

\begin{lstlisting}
maintainer@example.org ssh-ed25519 AAAA...TRUSTED_PUBLIC_KEY
\end{lstlisting}

Строка выше --- только описание формата; пользователь должен подставить
\emph{реальный} ранее проверенный публичный ключ, не копировать пример.

\begin{lstlisting}
bash scripts/verify_manifest_signature.sh \
    manifest.txt manifest.txt.sig trusted_signers maintainer@example.org
\end{lstlisting}

Скрипт проверяет домен \texttt{mangen-manifest} и доверенного
подписанта. Если манифест подменён, ключ не совпадает или подпись повреждена,
команда завершается с ошибкой. Для проверки содержимого каталога после
проверки подписи необходимо независимо пересчитать хэши его файлов;
подпись сама по себе доказывает неизменность списка, но не байтов файлов.

\section{Схема двухкоммитного релиза}

\begin{enumerate}
\item Коммит A содержит производственный код, тесты, \texttt{scripts/},
\texttt{README.md}, LaTeX-исходники и готовые PDF.
\item \texttt{scripts/package\_release.sh} создаёт ZIP из исходников
коммита A, добавляя \texttt{VERSION-COMMIT} с идентификатором A.
\item SHA-256 ZIP сохраняется в отдельный TXT в каталоге
\texttt{safeboard/}.
\item Коммит B добавляет только соответствующий TXT; сам ZIP не изменяется
после вычисления суммы и не входит в коммит B.
\end{enumerate}

Это исключает самоссылку: SHA-256 ZIP нельзя вписать в файл, который
уже находится внутри того же ZIP, не изменив архив. Встроенный
\texttt{VERSION-COMMIT} также создаётся на этапе упаковки: файл внутри
Git-коммита невозможно автоматически заполнить hash самого этого коммита.

\section{Ограничения}

Ссылка на файл может указывать за пределы дерева.
Файлы могут измениться между обнаружением и чтением; последовательные SHA
при изменении данных могут различаться. Защищённые каталоги приводят к
частичному результату. Для доказательной инвентаризации рекомендуется
неизменяемый набор данных, снимок файловой системы или изолированная копия.

\addcontentsline{toc}{section}{Список литературы}
\begin{thebibliography}{9}
\bibitem{ssh} OpenSSH. ssh-keygen(1), signature operations. \url{https://man.openbsd.org/ssh-keygen}
\bibitem{cmake} Kitware. CMake documentation. \url{https://cmake.org/cmake/help/latest/}
\bibitem{gnu} GNU Coreutils. SHA-2 utilities. \url{https://www.gnu.org/software/coreutils/manual/}
\bibitem{sha} NIST. FIPS 180-4. \url{https://csrc.nist.gov/pubs/fips/180-4/upd1/final}
\end{thebibliography}
\end{document}
